% ============================================================
% AgriTech SeedCounter – Classical CV Pipeline Report
% ============================================================
\documentclass[12pt,a4paper]{article}

% Packages
\usepackage[margin=2.5cm]{geometry}
\usepackage{graphicx}
\usepackage{booktabs}
\usepackage{multirow}
\usepackage{amsmath}
\usepackage{hyperref}
\usepackage{float}
\usepackage{subcaption}
\usepackage{xcolor}
\usepackage{caption}
\usepackage{array}
\usepackage{parskip}

\hypersetup{
    colorlinks=true,
    linkcolor=blue!70!black,
    urlcolor=blue!70!black,
    citecolor=blue!70!black
}

\title{
    \large Computer Vision Assignment 1\\[0.4em]
    \LARGE \textbf{Automated Seed Counting via Classical\\Image Processing Pipeline}\\[0.4em]
    \large AgriTech SeedCounter
}
\author{[Student Name]}
\date{\today}

\begin{document}

\maketitle
\tableofcontents
\newpage

% ============================================================
\section*{Abstract}
\addcontentsline{toc}{section}{Abstract}
% ============================================================

Accurate seed counting is a critical task in precision agriculture, enabling
farmers and researchers to assess sowing density, crop yield, and seed
viability. Manual counting is error-prone and infeasible at scale. This report
presents a classical image processing pipeline for automated seed counting,
comprising six sequential stages: preprocessing (noise reduction, CLAHE),
K-Means and DBSCAN clustering, adaptive thresholding, morphological
refinement, connected component analysis (CCA), and distance-transform
watershed for touching-seed separation. The pipeline was evaluated on a
dataset of 141 images spanning 1–144 seeds. On a representative 20-image
stratified sample, the baseline configuration achieved $50\%$ accuracy within
a $\pm 10\%$ error threshold (MAE\,=\,21.25, RMSE\,=\,34.23). A systematic
ablation study of seven configuration variants revealed that adaptive
thresholding consistently outperforms Otsu, and that morphological closing is
essential for high-density images. The primary failure mode is seed
under-segmentation at densities above \(\sim\!70\) seeds per image, where
merged foreground blobs cause watershed marker fragmentation.

\newpage

% ============================================================
\section{Introduction}
% ============================================================

\subsection{Problem Statement and Real-World Context}

Seed counting underlies several agricultural workflows: quality control in
seed production, sowing-rate calibration, and germination-rate studies.
Commercially, pharmaceutical and food companies also require accurate counting
of small, granular objects on trays. The traditional approach---human manual
counting---is both tedious (a tray of 140 seeds requires several minutes) and
subject to fatigue errors. Automated counting via computer vision offers a
scalable, consistent alternative.

The input images in this assignment were captured under controlled laboratory
lighting on a uniform background, with seeds arranged in clusters of varying
density (1 to 144 seeds per image, $3024 \times 3024$ pixels). Ground truth
counts are encoded in the image filename.

\subsection{Challenges in Automated Seed Counting}

Several factors make automated seed counting non-trivial even under controlled
conditions:

\begin{itemize}
    \item \textbf{Touching and overlapping seeds.} At high densities, adjacent
          seeds merge into a single foreground blob, making per-object
          segmentation difficult.
    \item \textbf{Reflections and shadows.} Seed surfaces produce specular
          reflections, and seed stacks cast shadows on their neighbours,
          creating intensity gradients within a single seed.
    \item \textbf{Scale variation.} The same pipeline must handle images with
          1 seed (very large object, most of the foreground) and 140+ seeds
          (small, densely packed objects).
    \item \textbf{Background noise.} Dust, scratches, and subtle texture in the
          tray surface can generate false positive foreground pixels.
\end{itemize}

\subsection{Objectives}

\begin{enumerate}
    \item Design and implement a six-stage classical CV pipeline for seed
          counting.
    \item Evaluate the impact of individual design choices (filter type,
          thresholding strategy, morphological configuration, segmentation
          source, and counting method) through a systematic ablation study.
    \item Identify the primary failure modes and propose directions for
          improvement.
\end{enumerate}

\newpage

% ============================================================
\section{Methodology}
% ============================================================

The pipeline is orchestrated in \texttt{main.py} and consists of the six
stages described below. Configuration is centralised in
\texttt{baseline\_code/config.yaml} for reproducibility.

\subsection{Stage 1: Preprocessing}

Raw BGR images are converted to grayscale before noise reduction. Three
filters are computed in parallel and automatically ranked by Laplacian
variance (a sharpness proxy):

\begin{itemize}
    \item \textbf{Gaussian blur} ($7 \times 7$ kernel, $\sigma = 0$): fast,
          effective for additive Gaussian noise.
    \item \textbf{Median filter} ($k = 5$): edge-preserving, better for salt-
          and-pepper noise.
    \item \textbf{Bilateral filter} ($d = 9$, $\sigma_c = 75$, $\sigma_s = 75$):
          smooths within regions while preserving sharp boundaries.
\end{itemize}

The filter producing the highest Laplacian variance is selected automatically.
In practice, Gaussian was selected on most images due to the relatively low
noise level. CLAHE (Contrast Limited Adaptive Histogram Equalization,
$\text{clipLimit}=2.0$, $8 \times 8$ grid) is always applied to the grayscale
image to improve local contrast before thresholding.

\begin{figure}[H]
    \centering
    \begin{subfigure}[b]{0.24\textwidth}
        \includegraphics[width=\textwidth]{../output/experiments/B2/100/grayscale_100.png}
        \caption{Grayscale Input}
    \end{subfigure}
    \begin{subfigure}[b]{0.24\textwidth}
        \includegraphics[width=\textwidth]{../output/experiments/B2/100/adaptive_100.png}
        \caption{Adaptive Threshold}
    \end{subfigure}
    \begin{subfigure}[b]{0.24\textwidth}
        \includegraphics[width=\textwidth]{../output/experiments/B2/100/morphological_100.png}
        \caption{Morphology}
    \end{subfigure}
    \begin{subfigure}[b]{0.24\textwidth}
        \includegraphics[width=\textwidth]{../output/experiments/B2/100/watershed_100.png}
        \caption{Watershed}
    \end{subfigure}
    \caption{Visualisation of the baseline pipeline stages (Image: 100.jpg).}
    \label{fig:pipeline_stages}
\end{figure}

\subsection{Stage 2: Clustering (K-Means and DBSCAN)}

Two unsupervised clustering algorithms are run on the preprocessed grayscale
image to produce alternative binary foreground masks:

\begin{itemize}
    \item \textbf{K-Means} ($K = 2$, pixel intensity feature): assigns each
          pixel to seed or background. The brighter cluster is taken as
          foreground.
    \item \textbf{DBSCAN} ($\varepsilon = 3$, $\text{min\_samples} = 10$, on
          pixel coordinates + intensity): identifies spatially coherent
          clusters. Cluster labels form a binary mask.
\end{itemize}

Clustering masks are compared against the thresholding mask via IoU for
logging. In the baseline configuration, the \emph{threshold mask} (Stage 3)
feeds downstream stages; K-Means and DBSCAN masks are explored as alternative
segmentation sources in Group D.

\subsection{Stage 3: Thresholding}

Two strategies are evaluated:

\begin{itemize}
    \item \textbf{Adaptive Gaussian thresholding} (block size\,=\,31,
          $C = 4$): computes a spatially varying threshold accounting for
          local illumination changes. This is the \emph{baseline} choice.
    \item \textbf{Otsu's global threshold}: finds a single optimal threshold
          by maximising inter-class variance. Effective on bimodal histograms
          but fails under uneven lighting.
\end{itemize}

\subsection{Stage 4: Morphological Refinement}

Binary masks are refined by a two-step morphological sequence applied to the
chosen segmentation mask:

\begin{enumerate}
    \item \textbf{Closing} (ellipse kernel $25 \times 25$, 2 iterations):
          fills intra-seed holes created by reflections and bridges narrow
          gaps between touching seeds.
    \item \textbf{Opening} (ellipse kernel $5 \times 5$, 1 iteration):
          removes small noise blobs and debris after closing.
\end{enumerate}

The order (Closing $\rightarrow$ Opening) is deliberate: closing first
consolidates seed regions, and opening then removes the newly introduced edge
artefacts.

\subsection{Stage 5: Object Counting --- CCA}

Connected Component Analysis (CCA, 8-connectivity) labels all foreground
blobs in the morphological mask. Components are filtered by:

\begin{itemize}
    \item \textbf{Area}: $[\text{min\_area},\, \text{max\_area}] = [8000,\, 200000]$
          pixels (full resolution).
    \item \textbf{Circularity}: $\frac{4\pi A}{P^2} \geq 0$ (threshold
          relaxed to zero in baseline; configurable).
    \item \textbf{Aspect ratio}: $\leq 3.0$ (excludes elongated debris).
    \item \textbf{Edge proximity}: components $>75\%$ outside the image
          boundary are excluded.
\end{itemize}

CCA is used as a cross-check and as the counting method in the CCA-only
ablation (E2).

\subsection{Stage 6: Watershed Segmentation (Post-processing)}

Watershed is the primary counting mechanism for touching seeds.

\begin{enumerate}
    \item \textbf{Distance transform} on the morphological mask produces a
          distance map where pixels far from the boundary have high values.
    \item \textbf{Sure foreground} thresholded at $0.4 \times$ peak distance
          identifies confident seed centres.
    \item \textbf{Sure background} obtained by dilating the mask with an
          $11 \times 11$ ellipse kernel ($\times 3$).
    \item \textbf{Marker labelling}: connected components of sure-foreground
          become markers; uncertain pixels (sure\_bg $-$ sure\_fg) are set to
          zero.
    \item \textbf{\texttt{cv2.watershed}} floods from each marker, drawing
          one-pixel boundaries between adjacent seeds.
    \item \textbf{Overlap correction}: watershed regions whose area exceeds
          \texttt{max\_area} are assumed to be merged seeds. Their estimated
          seed count is $\lceil A / \tilde{A} \rceil$ where $\tilde{A}$ is the
          median area of ``clean'' (area-bounded) watershed regions.
\end{enumerate}

A marker cap of 500 prevents pathological slowdowns on images where noise
fragments the sure-foreground into thousands of regions.

\subsection{Comparative Analysis of Alternative Approaches}

Seven variant configurations were evaluated against the baseline using a
20-image stratified sample and the experiment runner
(\texttt{run\_experiments.py}). Results are presented in Section~\ref{sec:results}.

\newpage

% ============================================================
\section{Results}
\label{sec:results}
% ============================================================

\subsection{Quantitative Results}

All experiments were evaluated on a 20-image stratified sample Covering the
full density range (1--140 seeds).

\begin{table}[H]
\centering
\caption{Pipeline variant comparison. Accuracy = fraction of images where
$|\hat{y} - y| / y \leq 10\%$. Best value per column is \textbf{bolded}.}
\label{tab:results}
\renewcommand{\arraystretch}{1.25}
\begin{tabular}{l l l r r r r}
\toprule
\textbf{ID} & \textbf{Group} & \textbf{Description} & \textbf{MAE} $\downarrow$ & \textbf{RMSE} $\downarrow$ & \textbf{Acc.\%} $\uparrow$ & \textbf{Fails} $\downarrow$ \\
\midrule
\textbf{★} & Baseline & Gaussian + Adaptive-G + Close→Open & \textbf{21.25} & \textbf{34.23} & \textbf{50.0} & \textbf{10} \\
\midrule
A2 & Filter    & Median filter (instead of Gaussian) & 21.60 & 34.68 & 45.0 & 11 \\
B2 & Threshold & Otsu global threshold               & 40.70 & 60.86 & 25.0 & 15 \\
C2 & Morphology & No closing (opening only)          & 66.80 & 81.00 & 0.0  & 20 \\
\midrule
D2 & Seg.\ Source & K-Means mask → CCA + Watershed  & 41.15 & 60.95 & 30.0 & 14 \\
D3 & Seg.\ Source & DBSCAN mask → CCA + Watershed   & 41.15 & 61.48 & 35.0 & 13 \\
\midrule
E2 & Counter  & CCA only (no watershed)              & 43.80 & 65.22 & 30.0 & 14 \\
E3 & Counter  & Blob detection only                  & 66.70 & 80.87 & 0.0  & 20 \\
\bottomrule
\end{tabular}
\end{table}

\subsection{Key Findings}

\paragraph{Adaptive thresholding is essential (B2 vs.\ baseline).}
Switching from adaptive to global Otsu nearly halves accuracy (25\% vs.\ 50\%).
Otsu computes a single threshold over the whole image; on images with lighting
gradients or seed shadows, this threshold is too low in dark regions (leaving
seeds undersegmented) and too high in bright regions (splitting seeds
prematurely).

\begin{figure}[H]
    \centering
    \begin{subfigure}[b]{0.48\textwidth}
        \includegraphics[width=\textwidth]{../output/experiments/B2/100/adaptive_100.png}
        \caption{Adaptive Thresholding (Baseline)}
    \end{subfigure}
    \hfill
    \begin{subfigure}[b]{0.48\textwidth}
        \includegraphics[width=\textwidth]{../output/experiments/B2/100/otsu_100.png}
        \caption{Otsu's Global Threshold (B2)}
    \end{subfigure}
    \caption{Comparison of thresholding methods on 100.jpg. Otsu fails to 
    segment individual seeds in dense regions compared to the adaptive approach.}
    \label{fig:threshold_comp}
\end{figure}

\paragraph{Morphological closing is critical (C2 vs.\ baseline).}
Removing the closing operation reduces accuracy to \textbf{0\%}—every image
exceeds the 10\% error threshold. Without closing, seed reflections create
holes in each seed's foreground mask. CCA then sees one large shell-like blob
per seed cluster rather than filled discs, making watershed unable to identify
distinct centres and drastically under-counting.

\paragraph{Watershed adds significant value over CCA alone (E2 vs.\ baseline).}
CCA-only (E2) achieves 30\% accuracy vs.\ 50\% for the full pipeline. For
touching seeds, CCA merges adjacent seeds into one component and counts it as
1; watershed's distance-transform markers split these blobs correctly.

\paragraph{Blob detection fails at high density (E3).}
Blob detection with LoG achieved 0\% accuracy. At high densities, the scale
parameter must be tuned per image, and overlapping blobs suppress each other
via the blob's non-maximum suppression.

\paragraph{Segmentation source matters less than expected (D2, D3 vs.\ baseline).}
K-Means and DBSCAN masks improve slightly over each other (35\% vs.\ 30\%)
but both fall short of the threshold-based baseline (50\%). The adaptive
threshold is computed on the locally contrast-enhanced (CLAHE) image, which
makes it more discriminative than raw-pixel K-Means clustering.

\subsection{Performance Across Image Conditions}

\begin{table}[H]
\centering
\caption{Baseline performance by image density category (from validation log).}
\label{tab:density}
\renewcommand{\arraystretch}{1.2}
\begin{tabular}{l r r r r}
\toprule
\textbf{Density} & \textbf{GT Range} & \textbf{Examples} & \textbf{Typical Error} & \textbf{Within 10\%} \\
\midrule
Sparse  & 1--10   & 1.jpg, 5.jpg, 10.jpg  & 0             & Yes \\
Medium  & 11--50  & 25.jpg, 50.jpg         & 0--11         & Yes (25.jpg), borderline (50.jpg) \\
Dense   & 51--144  & 100.jpg, 120.jpg+      & 90--130       & No \\
\bottomrule
\end{tabular}
\end{table}

The pipeline achieves near-perfect accuracy on sparse images ($\leq 25$
seeds). Performance degrades sharply above $\sim\!70$ seeds, where watershed
produces fewer than 20\% of the true markers due to blob merging.

\newpage

% ============================================================
\section{Error Analysis}
% ============================================================

\subsection{Primary Failure Mode: Blob Merging at High Density}

At densities above $\sim\!70$ seeds, the morphological mask consists of a
single connected blob covering nearly all seeds simultaneously. The distance
transform of this blob has a broad, flat peak rather than distinct local
maxima at each seed centre. The sure-foreground threshold (0.4 $\times$ peak)
then captures only one or two high-distance pixels, generating very few
markers—resulting in a severe under-count.

\textbf{Quantitative evidence:} On \texttt{100.jpg} (GT\,=\,100), CCA
returns 1 (the whole mass is one blob), watershed returns 10 (10 markers
found in the distance transform plateau). Error\,=\,90\%.

\begin{table}[H]
\centering
\caption{Selected per-image results from the validation log (baseline config).}
\renewcommand{\arraystretch}{1.2}
\begin{tabular}{l r r r r l}
\toprule
\textbf{Image} & \textbf{GT} & \textbf{CCA} & \textbf{Watershed} & \textbf{Error \%} & \textbf{Status} \\
\midrule
1.jpg   & 1   & 1  & 1  & 0.0\%  & \textcolor{green!60!black}{Correct} \\
5.jpg   & 5   & 5  & 5  & 0.0\%  & \textcolor{green!60!black}{Correct} \\
10.jpg  & 10  & 10 & 10 & 0.0\%  & \textcolor{green!60!black}{Correct} \\
25.jpg  & 25  & 23 & 25 & 0.0\%  & \textcolor{green!60!black}{Correct} \\
50.jpg  & 50  & 39 & 46 & 8.0\%  & \textcolor{orange}{Near-miss} \\
100.jpg & 100 & 1  & 10 & 90.0\% & \textcolor{red}{Failure} \\
\bottomrule
\end{tabular}
\end{table}

\subsection{Contributing Causes}

\begin{enumerate}
    \item \textbf{Seed touching at high density.} Seeds in the $50$--$144$
          range physically overlap in the image plane. Morphological closing
          then bridges the remaining gaps, creating one super-blob.
    \item \textbf{Uniform distance-transform plateau.} When seeds form a
          large compact blob, the distance transform has a wide plateau near
          the blob centre. The sure-foreground threshold ($0.4 \times$
          peak) captures only the very centre cluster of pixels, not one
          local maximum per seed.
    \item \textbf{Reflections breaking the mask.} Specular highlights on
          seed surfaces appear as dark pixels (after inversion) or white
          holes in the foreground mask. Despite closing, these can prevent
          some seeds from being properly filled, reducing the distance
          transform peak inside them.
    \item \textbf{Marker cap (500) and fallback.} For images where the
          sure-foreground fragments into $>500$ components (e.g.\ due to
          noise on the morphology mask), the watershed step is skipped and
          pure sure-fg CCA is used. This fallback under-counts when
          fragmentation reflects noise rather than seeds.
\end{enumerate}

\subsection{Failure Cases}

The most consistent failure pattern is observed on images with filenames
$\geq 70$ (GT $\geq 70$). The baseline achieves 0/10 correct predictions on
this subset of the 20-image sample, while achieving 10/10 on filenames $\leq
50$ (GT $\leq 50$). This clean density threshold suggests the failure
mechanism is structural (blob merging) rather than stochastic.

\newpage

% ============================================================
\section{Discussion}
% ============================================================

\subsection{Limitations of Classical Image Processing}

\paragraph{Fixed area thresholds.}
The \texttt{min\_area} and \texttt{max\_area} parameters in CCA and watershed
post-processing are set for full-resolution ($3024 \times 3024$) images. Any
downscaling for speed requires proportional scaling of these thresholds, adding
a brittle preprocessing dependency.

\paragraph{Non-adaptive distance threshold.}
The watershed sure-foreground threshold (0.4) is a single global parameter.
For a low-density image with well-separated seeds, 0.4 works well. For
a high-density image with merged blobs, the effective threshold should be
\emph{higher} to isolate individual seed centres. Adapting this parameter
per-image is a clear improvement direction.

\paragraph{Colour-agnostic pipeline.}
The entire pipeline operates on a single grayscale channel. Seed colour
differs noticeably from tray background colour. Incorporating a simple
colour-based mask (e.g.\ hue thresholding in HSV space) could provide a
cleaner initial foreground estimate than pure intensity thresholding.

\subsection{Scenarios Where the Approach Works Well}

\begin{itemize}
    \item \textbf{Sparse to medium images ($\leq 50$ seeds):} CCA perfectly
          separates all seeds; watershed correctly resolves the handful of
          touching pairs. Near-zero error.
    \item \textbf{Uniform lighting:} Adaptive thresholding is highly effective
          when block-level contrast is consistent, which holds for standard
          laboratory setups.
    \item \textbf{Well-separated seeds:} Both CCA and watershed are provably
          accurate when seeds do not touch (each seed $\rightarrow$ one
          component $\rightarrow$ one marker).
\end{itemize}

\subsection{Scenarios Where the Approach Fails}

\begin{itemize}
    \item \textbf{High density ($>70$ seeds):} Blob merging causes catastrophic
          under-counting regardless of thresholding or morphological strategy.
    \item \textbf{Extreme lighting:} Very strong directional lighting creates
          per-seed shadows that mimic separate objects, causing over-counting.
    \item \textbf{Seed pile-up (3D):} Seeds placed on top of each other violate
          the flat-tray assumption; their projected area is smaller than a
          single-layer seed.
\end{itemize}

\subsection{Ideas for Improvement}

\begin{itemize}
    \item \textbf{High-resolution local analysis:} Processing sub-windows of the
          image at higher resolution to resolve fine boundaries between seeds.
    \item \textbf{Adaptive watershed parameters:} Use image-level density
          estimates (e.g.\ foreground pixel fraction) to automatically tune
          the sure-foreground threshold on a per-image basis.
    \item \textbf{Multi-scale distance transform:} Apply the distance
          transform at multiple scales and aggregate markers to handle both
          small (sparse) and large (dense) seed configurations.
    \item \textbf{Shape-priors and Elliptical Fitting:} Leverage the fact that
          most seeds are approximately elliptical to fit parametric models to
          fragmented blobs, improving count accuracy in dense clusters.
\end{itemize}

\subsection{Real-World Deployment Considerations}

\begin{itemize}
    \item \textbf{Processing time:} The baseline pipeline takes approximately
          5--8 seconds per full-resolution image (including DBSCAN clustering).
          For a production tray-scanner running at 200 trays/hour, a GPU-based
          pipeline or image downscaling would be required.
    \item \textbf{Calibration:} The area thresholds and morphological kernel
          sizes are calibrated for the specific seed type and camera setup used
          here. Deployment on a different seed or camera requires recalibration.
    \item \textbf{Lighting control:} The pipeline's robustness can be greatly
          improved by controlling the imaging environment (telecentric lens,
          diffuse back-lit tray), reducing the need for adaptive thresholding.
\end{itemize}

\newpage

% ============================================================
\section{Conclusion}
% ============================================================

This report presented a six-stage classical image processing pipeline for
automated seed counting, evaluated through a structured ablation study. The
key findings are:

\begin{enumerate}
    \item \textbf{Adaptive thresholding} is the single most important
          algorithmic choice. Replacing it with Otsu reduces accuracy by
          25 percentage points on the 20-image sample.
    \item \textbf{Morphological closing} is indispensable. Removing it
          reduces accuracy to 0\% by creating holey masks that prevent
          watershed from finding meaningful markers.
    \item \textbf{Watershed} adds 20 percentage points over CCA alone,
          demonstrating the value of a touching-seed separation step.
    \item \textbf{Blob detection and DBSCAN segmentation} underperform
          the baseline for this dataset, confirming that adaptive thresholding
          on a CLAHE-enhanced image is the most effective foreground estimation
          strategy with classical methods.
    \item \textbf{High-density images remain the core unsolved challenge.}
          All classical variants fail on images with $>70$ seeds due to blob
          merging, with errors exceeding 80\% on the densest images.

\end{enumerate}

Future work should focus on density-adaptive watershed parameters as an
immediate improvement and on shape-aware object fitting as a longer-term 
solution that generalises across density levels without parameter re-tuning.

\end{document}
